% Generated from project outputs. Do not edit by hand.
\begin{table}[!htbp]
\centering
\begin{threeparttable}
\caption{Labor income by subsample: 2019 moments}
\label{tab:annual-2019-y1-moments}
\small
\setlength{\tabcolsep}{4pt}
\renewcommand{\arraystretch}{1.15}
\begin{tabularx}{\linewidth}{@{}Xrrrrr@{}}
\toprule
Sample & N & Mean & SD & Min & Max \\
\midrule
General & 94 & 2,730.7 & 1,581.7 & 342.7 & 6,230.1 \\
Men & 94 & 2,946.6 & 1,678.1 & 384.8 & 8,300.7 \\
Women & 94 & 2,368.1 & 1,704.9 & 253.6 & 7,656.7 \\
Urban & 94 & 2,929.5 & 1,617.4 & 308.7 & 6,376.1 \\
Rural & 94 & 2,315.6 & 1,514.2 & 278.7 & 9,510.9 \\
Indigenous & 94 & 2,133.2 & 1,431.6 & 237.4 & 7,537.6 \\
Non-Indigenous & 94 & 2,963.2 & 1,646.7 & 314.4 & 6,711.9 \\
Bottom 99 percent & 94 & 2,510.9 & 1,355.1 & 350.8 & 5,312.1 \\
Bottom 90 percent & 94 & 2,102.3 & 966.2 & 350.8 & 4,651.1 \\
Bottom 50 percent & 94 & 1,203.4 & 468.7 & 255.8 & 2,453.4 \\
\bottomrule
\end{tabularx}
\begin{tablenotes}[flushleft]
\footnotesize
\item Monthly quetzales. Unweighted distribution of survey-weighted cohort means before transition matching. N counts cohort cells, not individuals. SD is dispersion, not a standard error. Subsamples overlap. Cumulative income definitions and treatment of missing components follow the merging scripts.
\end{tablenotes}
\end{threeparttable}
\end{table}
